\section{函数矩阵对纯量的导数与积分}
\subsection{函数矩阵的基本运算}
%@see: 《矩阵分析（第3版）》（史荣昌、魏丰） P229
以实变量\(x\)的实值函数\(a_{ij}(x)\ (i=1,2,\dotsc,m;j=1,2,\dotsc,n)\)为元素的矩阵\begin{equation*}
	A(x)
	\defeq
	\begin{bmatrix}
		a_{11}(x) & a_{12}(x) & \dots & a_{1n}(x) \\
		a_{21}(x) & a_{22}(x) & \dots & a_{2n}(x) \\
		\vdots & \vdots & & \vdots \\
		a_{m1}(x) & a_{m2}(x) & \dots & a_{mn}(x) \\
	\end{bmatrix}
\end{equation*}
称为\DefineConcept{函数矩阵}.

当行数\(m=1\)时，\(A(x)\)称为\DefineConcept{函数列向量}.
当列数\(n=1\)时，\(A(x)\)称为\DefineConcept{函数行向量}.

函数矩阵的加法、数量乘法，两个函数矩阵的乘法，函数矩阵的转置，
都与之前定义的常数矩阵的相应运算完全相同.
类似地，可以定义函数矩阵的伴随矩阵、逆矩阵.
\begin{definition}
%@see: 《矩阵分析（第3版）》（史荣昌、魏丰） P230 定义7.1.1
设\(A(x)\)是\(n\)阶函数矩阵.
若存在\(n\)阶函数矩阵\(B(x)\)
使得对于任何\(x \in [a,b]\)
都有\begin{equation*}
	A(x) B(x) = B(x) A(x) = E,
\end{equation*}
其中\(E\)是\(n\)阶单位矩阵，
则称
“\(A(x)\)是区间\([a,b]\)上的一个\DefineConcept{可逆（函数）矩阵}”
或称“\(A(x)\)在区间\([a,b]\)上\DefineConcept{可逆}”；
并称“\(B(x)\)是\(A(x)\)的\DefineConcept{逆（函数）矩阵}”，
记\(A^{-1}(x) \defeq B(x)\).
\end{definition}

\begin{theorem}
%@see: 《矩阵分析（第3版）》（史荣昌、魏丰） P230 定理7.1.1
\(n\)阶函数矩阵\(A(x)\)在区间\([a,b]\)上可逆的充分必要条件是
\(\DeterminantA{A(x)}\)在\([a,b]\)上处处不等于零，
并且\begin{equation*}
	A^{-1}(x)
	= \frac{1}{\DeterminantA{A(x)}}
	A^*(x),
\end{equation*}
其中\(A^*(x)\)是\(A(x)\)的伴随矩阵.
\end{theorem}
